\chapter{แบบจำลองปัญญาประดิษฐ์ประมวลผลบนขอบ}
\label{chap:edge_ml}

\section{สถาปัตยกรรมปัญญาประดิษฐ์บนขอบ (Pure Dart Edge ML Architecture)}
จุดเด่นสำคัญที่สุดประการหนึ่งของโครงการ SoilpHTxAI คือ การพัฒนาแบบจำลองปัญญาประดิษฐ์ที่ทำงานบนตัวเครื่องสมาร์ทโฟนโดยตรง (On-Device Edge ML) ทั้งในขั้นตอนการอนุมานผล (Inference) และขั้นตอนการฝึกฝนแบบจำลองใหม่ (In-App Model Re-training) ด้วยภาษา Dart บริสุทธิ์ 100\% โดยปราศจากการพึ่งพาไลบรารีภายนอกที่เป็นภาระ เช่น TensorFlow, PyTorch หรือเซิร์ฟเวอร์คลาวด์ภายนอก

แนวทางนี้ช่วยปลดล็อกข้อจำกัดของการเกษตรในพื้นที่ห่างไกล ซึ่งมักไม่มีสัญญาณอินเทอร์เน็ตหรือสัญญาณไม่เสถียร ทำให้เกษตรกรและนักวิจัยสามารถปรับจูนเครื่องมือวัดเข้ากับสารละลายบัฟเฟอร์มาตรฐานในห้องปฏิบัติการภาคสนามได้ทันที

\section{ฟังก์ชันการแมปปิ้งพหุนามที่ไม่เป็นเชิงเส้น}
แบบจำลองใช้สมการการถดถอยพหุนามหลายตัวแปร (Multivariate Polynomial Regression) เพื่อสร้างความสัมพันธ์แบบผกผันระหว่างค่าศักย์ไฟฟ้าที่วัดได้ ($E$ ในหน่วย mV) และอุณหภูมิ ($T$ ในหน่วย $^\circ\text{C}$) สู่ค่าความเป็นกรด-ด่างจริง ($pH_{pred}$) ดังนี้

\begin{formulabox}{สมการแบบจำลองปัญญาประดิษฐ์ชดเชยความคลาดเคลื่อน}
\begin{equation}
\text{pH}_{pred} = w_0 + w_1 E + w_2 T + w_3 (E \cdot T) + w_4 E^2 + w_5 T^2
\label{eq:poly_model}
\end{equation}
โดยที่
\begin{itemize}
    \item $w_0$ คือ ค่าคงที่ไบแอสของระบบ (Bias Constant) ชดเชยจุดสมมูล Isopotential
    \item $w_1, w_2$ คือ สัมประสิทธิ์เชิงเส้นของศักย์ไฟฟ้าและอุณหภูมิ
    \item $w_3$ คือ สัมประสิทธิ์อันตรกิริยาคู่ควบข้ามตัวแปร (Cross-coupling Interaction) ระหว่าง $E$ และ $T$
    \item $w_4, w_5$ คือ สัมประสิทธิ์กำลังสอง ชดเชยความโค้งไม่เป็นเชิงเส้น (Non-linear Curvature) และความคลาดเคลื่อนของเยื่อแก้ว
\end{itemize}
\end{formulabox}

\section{การแก้ระบบสมการปรกติด้วยอัลกอริทึม Gaussian Elimination}
ในการฝึกฝนแบบจำลองบนตัวเครื่องจากจุดข้อมูลสอบเทียบ $n$ จุด ระบบสร้างเมทริกซ์คุณลักษณะ $X$ ขนาด $n \times 6$ ซึ่งประกอบด้วยเวกเตอร์ของแต่ละจุดข้อมูล $x_i = [1.0, E_i, T_i, E_i T_i, E_i^2, T_i^2]$ จากนั้นระบบจะสร้างระบบสมการปรกติ (Normal Equation) ของวิธีกำลังสองน้อยที่สุด (Ordinary Least Squares หรือ OLS):

\begin{equation}
(X^T X) \cdot W = X^T \cdot Y
\label{eq:normal_equation}
\end{equation}

ให้ $A = X^T X$ เป็นเมทริกซ์ขนาด $6 \times 6$ และ $B = X^T Y$ เป็นเวกเตอร์ขนาด $6 \times 1$ ระบบ Dart แก้ระบบสมการ $A \cdot W = B$ ด้วยวิธีขจัดแบบเกาส์เซียนพร้อมการเลือกสมาชิกหลักบางส่วน (Gaussian Elimination with Partial Pivoting) ร่วมกับการใส่ค่าควบคุมความราบเรียบแบบริดจ์ (Ridge Regularization $\lambda = 10^{-6}$) เพื่อป้องกันปัญหาเมทริกซ์เอกฐาน (Singular Matrix)

\begin{table}[H]
\caption{ผลการประเมินสมรรถนะของแบบจำลอง Edge ML ที่ฝึกฝนบนสมาร์ทโฟน}
\label{tab:edge_ml_metrics}
\centering
\begin{tabularx}{\textwidth}{p{4.5cm}cZZ}
\toprule
\textbf{ตัวชี้วัดสมรรถนะ} & \textbf{สัญลักษณ์} & \textbf{ค่าที่ได้บนเครื่อง} & \textbf{เกณฑ์การยอมรับ} \\
\midrule
สัมประสิทธิ์การตัดสินใจ & $R^2$ & 0.9998 & $> 0.990$ (ดีเยี่ยม) \\
รากที่สองของค่าคลาดเคลื่อนกำลังสองเฉลี่ย & RMSE & 0.0152 pH & $< 0.050$ pH \\
ค่าคลาดเคลื่อนสัมบูรณ์เฉลี่ย & MAE & 0.0121 pH & $< 0.040$ pH \\
เวลาในการฝึกฝนและคำนวณ & Duration & 118 ms & $< 500$ ms (ประมวลผลทันที) \\
จำนวนตัวอย่างที่ใช้ฝึกฝน & Sample Count & 118 จุด (70\%) & ครอบคลุมทุกสภาวะ \\
\bottomrule
\end{tabularx}
\end{table}

\section{ระบบสลับเปลี่ยนโมเดลแบบ Dynamic Hot-Swapping}
แอปพลิเคชันรองรับการสลับโมเดลได้แบบเรียลไทม์ โดยผู้ใช้สามารถกดเปิดหน้าต่าง Edge ML Trainer จากหน้าชุดข้อมูล จากนั้นกดปุ่มเพื่อฝึกฝนโมเดล เมื่อได้ผลการฝึกฝนที่พึงพอใจ ผู้ใช้สามารถกดปุ่ม ``นำไปใช้งานทันที'' เพื่อฉีดค่าน้ำหนักสัมประสิทธิ์ชุดใหม่ ($w_0 - w_5$) เข้าสู่ระบบการอนุมานผลของเซนเซอร์สดได้โดยไม่ต้องรีสตาร์ตแอปพลิเคชัน พร้อมกันนี้ยังมีปุ่ม ``รีเซ็ตสู่ค่าเริ่มต้น'' หากต้องการคืนค่ากลับสู่ค่าน้ำหนักมาตรฐานการวิจัย RBRU Baseline 2569

\clearpage
